\documentclass[10pt,a4paper]{amsart}
\usepackage[margin=19mm]{geometry}
\usepackage{amsmath,amssymb,mathtools}
\usepackage[scr=rsfs,cal=euler]{mathalfa}
\usepackage{xfrac}
\usepackage[unicode,hidelinks,bookmarks=false]{hyperref}
\newcommand{\ie}{i.e.\ }
\newcommand{\cf}{cf.\ }
\newcommand{\resp}{respectively\ }
\newcommand{\Subsection}{\subsection}
\providecommand{\nobreakdash}{\nobreak\hbox{-}}
\setlength{\emergencystretch}{2em}
\multlinegap0pt
\usepackage{bm}
\usepackage{bbm}
\usepackage{dsfont}
\usepackage{xspace}
\usepackage{cancel}
\usepackage{mathtools}
\usepackage{amsthm}
\makeatletter
\@ifundefined{theorem}{\newtheorem{theorem}{Theorem}}{}
\@ifundefined{corollary}{\newtheorem{corollary}[theorem]{Corollary}}{}
\@ifundefined{proposition}{\newtheorem{proposition}[theorem]{Proposition}}{}
\makeatother
\newcommand{\cW}{{\mathcal{W}}}
\newcommand{\sfA}{\mathsf{A}}
\renewcommand{\Cap}{\operatorname{Cap}}
\title[Weighted residual polynomials on a circular arc]{Weighted residual polynomials on a circular arc}
\author{Jacob S. Christiansen, Benjamin Eichinger, Olof Rubin, Maxim Zinchenko}
\date{Source: arXiv:2602.05428v1 (CC BY 4.0)}
\begin{document}
\maketitle

\noindent\emph{Source and licence.} Weighted residual polynomials on a circular arc. Version arXiv:2602.05428v1 (CC BY 4.0), distributed under the Creative Commons Attribution 4.0 International licence, as recorded in the arXiv licence metadata for this exact version and shown on the abstract page https://arxiv.org/abs/2602.05428v1. Full rights notice: licenses/arxiv-2602.05428-CC-BY-4.0.txt.\par\smallskip

\noindent\textit{Editorial scope.} Counted unit: the Corollary whose source label is \texttt{None} in the article (source0.tex lines 1798--1804, source label \texttt{None}), delivered together with the article's own complete original proof of it (source0.tex lines 1806--1825, one \texttt{proof} environment of 20 lines). No mathematics is written, repaired or re-derived: statement and proof are the article's own text line for line. Required context retained uncounted: eq:weight\_function\_form (lines 1240--1243); hw (lines 1688--1691); eq:outer\_function (lines 1694--1697); eq:residual\_at\_zero\_limit (lines 1703--1708); ca (lines 1763--1767); eq:residual\_at\_infinity\_limit (lines 1776--1785). Every definition, display and label of those lines is retained unchanged. Omitted from this package and not redistributed: 1--1239, 1244--1687, 1692--1693, 1698--1702, 1709--1762, 1768--1775, 1786--1797, 1805--1805, 1826--3164. Nothing inside a retained excerpt is omitted or altered: every retained body is delivered exactly as the article's lines are written, so no omission receipt of any kind accompanies this package. Citations: the retained excerpts carry no citation command at all, so no bibliography entry of the article is transcribed and no reference is redistributed. Reference adaptation: none was needed and none was made; every reference inside the delivered unit resolves inside it, so no unresolved reference or citation is produced. Printed slips kept literal: no printed word, symbol, hypothesis or number of the counted statement or of its proof was altered. Layout and class: the article's own class is \texttt{amsart} with packages including inputenc, fontenc, hyperref, amsmath, amsthm, amsfonts, amssymb, amsbsy, subcaption, pdfpages, fancyhdr, geometry, setspace, xcolor; the wrapper is the lane's pinned \texttt{amsart} editorial wrapper at 10pt on A4, which restates the theorem-like environments the retained text needs and adds the article's own macro definitions that the retained text needs (\texttt{\textbackslash Cap}, \texttt{\textbackslash cW}, \texttt{\textbackslash sfA}), with extra wrapper packages (bm, bbm, dsfont, xspace, cancel, mathtools). The delivered numbering is the unit's own. Assets: no retained span contains a figure environment, a graphics-inclusion command, a PDF-page inclusion, a table or a TikZ picture; no image file is copied into this package, the package declares no asset dependency and no image file is redistributed with it. Licence: version arXiv:2602.05428v1 of this article is distributed under the Creative Commons Attribution 4.0 International licence, as recorded in the arXiv licence metadata for this exact version and shown on the abstract page https://arxiv.org/abs/2602.05428v1. A full rights notice is delivered at licenses/arxiv-2602.05428-CC-BY-4.0.txt. Characters: every retained excerpt is pure ASCII, so the delivered engine, XeLaTeX with its default Latin Modern text font, reports no missing character.
\begin{equation}
	w(u) = e^{i\beta}\Biggl[\,\prod_{j=1}^{k}(u-u_j)\Biggr]^{-1},
	\label{eq:weight_function_form}
\end{equation}

\begin{equation}
\label{hw}
   h(u) := -\int_{\Gamma_\alpha}\log |w(x)|\,\omega(u,dx,\Omega_{\alpha}).
\end{equation}

\begin{equation}
\label{eq:outer_function}
   F_w(u, u_0) := \exp \Bigl[ h(u)+i \widetilde{h}(u) \Bigr]
\end{equation}

\begin{equation}
\label{eq:residual_at_zero_limit}
   b_{\Omega_\alpha}(u,\infty)^n R_{n}^{\Gamma_\alpha,w}(u,0) \longrightarrow  
   \frac{u+1+\sqrt{(u-e^{i\alpha})(u-e^{-i\alpha})}}{2(u+1)}
   \frac{b_{\Omega_\alpha}(u, -1)}{b_{\Omega_\alpha}(u,\infty)} F_w(u, 0),
\end{equation}

\begin{equation}
\label{ca}
   b_{\Omega_\alpha}(\infty, 0) = \lim_{u\to\infty} ub_{\Omega_\alpha}(u, \infty) =  
   \Cap(\Gamma_\alpha)=\sin(\alpha/2).
\end{equation}    

\begin{proposition}
Let $w$ be the weight function defined in \eqref{eq:weight_function_form}. Then \eqref{eq:residual_at_zero_limit} holds, and with the normalization $b_{\Omega_\alpha}(\infty,-1)>0$, we have
\begin{equation}
	b_{\Omega_{\alpha}}(u,\infty)^nR_{n}^{\Gamma_\alpha,w}(u,\infty)\longrightarrow 
	\frac{u+1-\sqrt{(u-e^{i\alpha})(u-e^{-i\alpha})}}{2(u+1)}
	\frac{b_{\Omega_\alpha}(u,-1)}{b_{\Omega_\alpha}(u,0)}F_w(u,\infty)
	\label{eq:residual_at_infinity_limit}
\end{equation}
uniformly on compact subsets of $\Omega_\alpha$ as $n\rightarrow \infty$.
\end{proposition}

\begin{corollary}
Let $w$ be the weight function defined in \eqref{eq:weight_function_form}. Then 
\begin{equation}
   \lim_{n\rightarrow \infty}\cW_n(\Gamma_\alpha,w) = 
   2\cos(\alpha/4)^2\exp\left(\int_{\Gamma_\alpha} \log |w(x)|\,d\mu_{\Gamma_\alpha}(x)\right).
\end{equation}
\end{corollary}

\begin{proof}
It suffices to compute the limit of the right-hand side of \eqref{eq:residual_at_infinity_limit} as $u\to\infty$. Since $s(z_\infty)=-1$, the first factor tends to $1$. 
To evaluate $b_{\Omega_\alpha}(\infty, -1)$, we note that a direct calculation using the Joukowski transformation gives
\[
   b_{\Omega_0}(z,0)  = \frac{-1+\sqrt{1-z^2}}{iz} = \frac{iz}{1+\sqrt{1-z^2}},
\]
where the square root is taken with a branch cut along $\sfA_0$ and chosen so that $\sqrt{1}=1$.
Substituting $z_\infty = -i\tan(\alpha/2)$, we obtain
\[
   b_{\Omega_\alpha}(\infty, -1) = b_{\Omega_0}(z_\infty, 0) 
   = \frac{\tan(\alpha/2)}{1+\sqrt{1+\tan^2(\alpha/2)}} % = \frac{\sin(\alpha/2)}{\cos(\alpha/2)+1} 
   = \frac{\sin(\alpha/2)}{2\cos^2(\alpha/4)}.
\]
However, the numerator cancels out, as $b_{\Omega_\alpha}(\infty, 0) = \sin(\alpha/2)$ by \eqref{ca}.
Finally, from \eqref{hw}--\eqref{eq:outer_function}, we see that 
\[
   F_w(\infty,\infty)=\exp\left(-\int_{\Gamma_\alpha} \log |w(x)|\,d\mu_{\Gamma_\alpha}(x)\right),
\]
since the normalization ensures that $F_w(\infty,\infty)>0$. This completes the proof.
\end{proof}

\end{document}
